\section{Weyl's theorem}

I use my proof in a former note use short exact sequence.

\begin{lemma}
    Let $L$ be a semisimple Lie algebra and $\phi : L \to \mathfrak{gl}(V)$ a representation. Then $\phi(L) \subseteq \mathfrak{sl}(V)$. In particular, $L$ acts trivially on any one-dimensional $L$-module.
\end{lemma}

\begin{proof}
    Since $L$ is semisimple, we have $L = [L,L]$. Hence
    \[
        \phi(L) = \phi([L,L]) = [\phi(L), \phi(L)].
    \]
    But $[\mathfrak{gl}(V), \mathfrak{gl}(V)] = \mathfrak{sl}(V)$, so $\phi(L) \subseteq \mathfrak{sl}(V)$.

    Let $W$ be a one-dimensional $L$-module, and let $\psi : L \to \mathfrak{gl}(W) \cong \Bbbk$. Then for all $x \in L$, we have $\operatorname{Tr}(\psi(x)) = 0$. Hence $\psi(x) = 0$, so $L$ acts trivially on $W$.
\end{proof}

\begin{theorem}
    Let $L$ be a semisimple Lie algebra and $\phi : L \to \mathfrak{gl}(V)$ a representation, where $V$ is finite-dimensional. Then $\phi$ is completely reducible. In other words, if $W$ is an $L$-submodule of $V$, then there exists a submodule $W'$ such that
    \[
        V = W \oplus W'.
    \]
\end{theorem}

\begin{remark}
    If $V$ is irreducible, then there is nothing to prove. Hence we may assume that $V$ is reducible.
\end{remark}

\begin{remark}
    Being completely reducible is equivalent to the following: for any submodule $W \subseteq V$, there exists a submodule $W'$ such that $V = W \oplus W'$. Furthermore, this is equivalent to saying that the short exact sequence
    \[
        0 \longrightarrow W \longrightarrow V \longrightarrow V/W \longrightarrow 0
    \]
    splits.
\end{remark}

\begin{definition}[Definition 2.17]
    The short exact sequence of a direct sum is called a \emph{splitting short exact sequence}, i.e.
    \[
        0 \longrightarrow N \xrightarrow{f} N \oplus P \xrightarrow{g} P \longrightarrow 0
    \]
    where $f : N \to N \oplus P$ is given by $n \mapsto (n,0)$, and
    $g : N \oplus P \to P$ is given by $(n,p) \mapsto p$.
\end{definition}

\begin{definition}[Definition 2.18]
    A short exact sequence
    \[
        0 \longrightarrow N \xrightarrow{u} M \xrightarrow{v} P \longrightarrow 0
    \]
    is called a \emph{splitting short exact sequence} if there exists an $R$-module isomorphism
    \[
        \theta : M \to N \oplus P
    \]
    such that the following diagram commutes:
    \[
        \begin{tikzcd}
            0 \arrow[r] & N \arrow[r, "u"] \arrow[d, equal]
            & M \arrow[r, "v"] \arrow[d, "\theta"]
            & P \arrow[r] \arrow[d, equal] & 0 \\
            0 \arrow[r] & N \arrow[r, "f"]
            & N \oplus P \arrow[r, "g"]
            & P \arrow[r] & 0
        \end{tikzcd}
    \]
    An intuition is that $u$ and $v$ behave like the standard embedding $f$ and projection $g$, respectively.
\end{definition}

\begin{theorem}[Theorem 2.9]
    Consider the following short exact sequence:
    \[
        0 \longrightarrow N \xrightarrow{u} M \xrightarrow{v} P \longrightarrow 0
    \]
    Then the following are equivalent:
    \begin{enumerate}
        \item There exists $u' : M \to N$ such that $u' \circ u = \mathrm{id}_N$.
        \item There exists $v' : P \to M$ such that $v \circ v' = \mathrm{id}_P$.
        \item There exists an isomorphism $\theta : M \to N \oplus P$ such that the diagram
              \[
                  \begin{tikzcd}
                      0 \arrow[r] & N \arrow[r, "u"] \arrow[d, equal]
                      & M \arrow[r, "v"] \arrow[d, "\theta"]
                      & P \arrow[r] \arrow[d, equal] & 0 \\
                      0 \arrow[r] & N \arrow[r, "f"]
                      & N \oplus P \arrow[r, "g"]
                      & P \arrow[r] & 0
                  \end{tikzcd}
              \]
              commutes.
    \end{enumerate}
    In other words, the short exact sequence splits.
\end{theorem}

\begin{proof}
    (1) $\Rightarrow$ (3): Define $\theta : M \to N \oplus P$ by
    \[
        \theta(m) = (u'(m), v(m)).
    \]
    Then the diagram commutes. We show $\theta$ is an isomorphism.

    \textbf{Injective:} Suppose $\theta(m) = (0,0)$. Then $v(m)=0$, so $m \in \operatorname{Im}(u)$, i.e. $m = u(n)$ for some $n \in N$. Hence
    \[
        0 = u'(m) = u'(u(n)) = n,
    \]
    so $m = u(n) = 0$.

    \textbf{Surjective:} Let $(n,p) \in N \oplus P$. Since $v$ is surjective, there exists $m \in M$ such that $v(m)=p$. Then
    \[
        v(m + u(x)) = v(m) = p \quad \text{for all } x \in N.
    \]
    Take $x = n - u'(m)$. Then
    \[
        u'(m + u(x)) = u'(m) + x = n,
    \]
    so
    \[
        \theta(m + u(x)) = (n,p).
    \]
    Thus $\theta$ is surjective.

    (3) $\Rightarrow$ (1): Since $\theta$ is an isomorphism and the diagram commutes, write
    \[
        \theta(m) = (u'(m), v(m)).
    \]
    Then for $n \in N$,
    \[
        \theta(u(n)) = (n,0),
    \]
    so $u'(u(n)) = n$, hence $u' \circ u = \mathrm{id}_N$.

    (3) $\Rightarrow$ (2) and (2) $\Rightarrow$ (3) are similar.
\end{proof}

Back to the proof of Weyl Theorem:

\begin{proof}
    What we want to show is $\forall W \subset V$ as submodule, the short exact sequence
    \[
        0 \longrightarrow W \longrightarrow V \longrightarrow V/W \longrightarrow 0
    \]
    splits. We begin with very special case, and step by step, we will prove the general case.

    \vspace{1em}
    \noindent \textbf{Step 1:} At first we prove if $W$ is irreducible with codimension 1 in $V$, then the sequence splits. i.e. $W$ has a direct sum complement in $V$.

    Let $\phi : L \to \mathfrak{gl}(V)$. Take the casimir element $C_\phi$ as before. Since $C_\phi$ commutes with $\phi(L)$, $\ker C_\phi$ is an $L$-submodule of $V$.

    Since $W$ is an $L$-submodule, we have $C_\phi(W) \subseteq W$.

    Since $W$ is of codimension 1, $L$ acts trivially on $V/W$, we have $C_\phi$ acts trivially on $V/W$. Therefore $C_\phi|_{V/W} = 0$.

    By Schur lemma, $C_\phi|_W = \lambda \cdot 1_W$. i.e. $\mathrm{Tr}(C_\phi|_W) = \lambda \cdot \dim W$. Recall $\mathrm{Tr}\, C_\phi = \dim L$ and $C_\phi|_{V/W} = 0$, we have $\lambda = \dim L / \dim W > 0$. Thus $C_\phi$ acts on $W$ as $\mathrm{diag}\{\lambda, \dots, \lambda\}$ where $\lambda > 0$. Thus $C_\phi(W) = W$, $C_\phi|_{V/W} = 0$. Take $\frac{1}{\lambda} \ker C_\phi$ then $\frac{1}{\lambda} \ker C_\phi$ is an $L$-submodule.

    Now consider $\forall v \in V$, $v = v - \frac{1}{\lambda}C_\phi(v) + \frac{1}{\lambda}C_\phi(v)$. Remark, since $C_\phi$ acts trivially on $V/W$ we have $C_\phi(V) \subseteq W$. Since $C_\phi\left(v - \frac{1}{\lambda}C_\phi(v)\right) = C_\phi(v) - \frac{1}{\lambda}C_\phi(C_\phi(v)) = C_\phi(v) - \frac{1}{\lambda} \cdot \lambda C_\phi(v) = 0$, $v - \frac{1}{\lambda}C_\phi(v) \in \ker C_\phi$. Thus $V = \frac{1}{\lambda}\ker C_\phi + \mathrm{im}\, C_\phi = \frac{1}{\lambda}\ker C_\phi + W$.

    Finally, if $\exists v_0$ such that $v_0 \in \ker C_\phi \cap \mathrm{im}\, C_\phi$. Then $C_\phi(v_0) = \lambda v_0 = 0$ shows $v_0 = 0$. So $V = \frac{1}{\lambda}\ker C_\phi \oplus W$.

    \vspace{1em}
    \noindent \textbf{Step 2:} If $W$ is not irreducible. Then we can find a non-trivial irreducible submodule $W_1$ such that $0 \neq W_1 \subset W$. Take
    \[
        0 \longrightarrow W/W_1 \longrightarrow V/W_1 \longrightarrow \frac{V/W_1}{W/W_1} \longrightarrow 0
    \]
    Since $\mathrm{codim}_V W = 1$, we have $\mathrm{codim}_{V/W_1} W/W_1 = 1$. And we can induction on $\dim W$.

    As $0 \neq W_1 \subset W$, $\dim W/W_1 < \dim W$. We can assume
    \[
        0 \longrightarrow W/W_1 \longrightarrow V/W_1 \longrightarrow \frac{V/W_1}{W/W_1} \longrightarrow 0
    \]
    splits. i.e. $V/W_1 = W/W_1 \oplus \tilde{W}/W_1$ for some $\tilde{W}$ and $\dim \tilde{W}/W_1 = 1$.

    Now consider
    \[
        0 \longrightarrow W_1 \longrightarrow \tilde{W} \longrightarrow \tilde{W}/W_1 \longrightarrow 0
    \]
    Since $\dim W_1 < \dim W$, by induction $\exists X$ such that $\tilde{W} = W_1 \oplus X$ where $\dim X = 1$. Write $X = \mathrm{span}\{x\}$. $\forall v \in V, v + W_1 = w + W_1 + \tilde{w} + W_1 = w + W_1 + w_1 + kx + W_1 = w + kx + W_1 = kx + W$ where $k \in F$. In other words, $V = W \oplus X$ this is a direct sum since $W, X \subset V$ and $W \cap X = 0$.

    \vspace{1em}
    \noindent \textbf{Step 3:} Take $V$ as an arbitrary $L$-module (reducible) and $W \subset V$ non-trivial submodule. Consider $\Hom_F(V, W)$. Remark, $\Hom_F(V, W)$ is the set of linear maps but not $L$-module homomorphisms. $\Hom_F(V, W)$ itself is an $L$-module.

    Define $\mathcal{V} \subset \Hom_F(V, W)$, $\mathcal{V} = \{f \in \Hom_F(V, W) \mid f|_W = a \cdot 1_W, a \in F\}$. Also define $\mathcal{W} = \{g \in \Hom_F(V, W) \mid g|_W = 0\}$. Then $\mathcal{V}, \mathcal{W}$ are $L$-modules and $\mathcal{W} \subset \mathcal{V}$.

    Actually, $\forall f \in \mathcal{V}, v \in V, x \in L$, $(x \cdot f)(v) = x \cdot f(v) - f(x \cdot v) = x \cdot (av) - a(x \cdot v) = 0$, i.e. $x \cdot (f(v)) = f(x \cdot v)$. So $f$ is an $L$-module homomorphism. Besides, it is obvious that $\mathrm{codim}_{\mathcal{V}} \mathcal{W} = 1$.

    So we have
    \[
        0 \longrightarrow \mathcal{W} \longrightarrow \mathcal{V} \longrightarrow \mathcal{V}/\mathcal{W} \longrightarrow 0
    \]
    splits. i.e. $\mathcal{V} = \mathcal{W} \oplus \tilde{\mathcal{V}}$ for some $\tilde{\mathcal{V}}$ and $\tilde{\mathcal{V}} = \mathrm{span}\{f\}$ for some $f \in \Hom_F(V, W)$ where $f|_W = a \cdot 1_W$ for some $0 \neq a \in F$.

    Now take $\frac{1}{a}f$ consider the origin short exact sequence
    \[
        \begin{tikzcd}
            0 \arrow[r] & W \arrow[r, "i", shift left=0.6ex] & V \arrow[l, "\frac{1}{a}f", dashed, shift left=0.6ex] \arrow[r, "\pi"] & V/W \arrow[r] & 0
        \end{tikzcd}
    \]
    $\frac{1}{a}f \circ i = \mathrm{id}_W$ shows this short exact sequence splits.

\end{proof}



